\documentclass[12pt]{article}

\usepackage[margin=1in]{geometry}
\usepackage[T1]{fontenc}
\usepackage[utf8]{inputenc}
\usepackage{lmodern}
\usepackage{setspace}
\usepackage{amsmath,amssymb,amsthm}
\usepackage{mathtools}
\usepackage[round]{natbib}
\usepackage{booktabs}
\usepackage{graphicx}
\usepackage{caption}
\usepackage{xcolor}
\usepackage{enumitem}
\usepackage{xurl}
\usepackage{xr-hyper}
\usepackage[colorlinks=true,linkcolor=blue!60!black,citecolor=blue!60!black,urlcolor=blue!60!black]{hyperref}
\usepackage[capitalise,nameinlink]{cleveref}

% References to the working paper carry the prefix M- and are resolved from
% main.aux, so main.tex must be compiled first.
\externaldocument[M-]{main}

\onehalfspacing
\graphicspath{{figures/}}
\captionsetup{font=small,labelfont=bf}
\setlist{itemsep=2pt,topsep=4pt}
\renewcommand{\thesection}{S\arabic{section}}
\numberwithin{equation}{section}
\crefname{section}{Section}{Sections}

% Notation
\newcommand{\bone}{\mathbf{1}}
\newcommand{\TB}{\mathrm{TB}}
\newcommand{\ToT}{\mathrm{ToT}}

\title{\textbf{Online Supplement to\\ ``Understanding Multilateral Tariffs and Exchange Rates''}}
\author{Jason Lu \and Dimitre Milkov}
\date{\vspace{-1.5em}}

\begin{document}

\maketitle

\noindent This supplement collects derivations, numerical results and data documentation that are referenced in, but not required for, the working paper. Section, equation, proposition, figure and table numbers without the prefix S refer to the working paper. Notation and the sign convention are those of the paper: a rise in $e_{Aj}$ or $\nu_j$ is a depreciation of $A$ against $j$. Numerical values use the US-like parameters $(\alpha_D,\alpha_T,\eta)=(0.8,0.4,1.5)$, $\rho_D=1.32$, unless stated otherwise.

\section{Three-country illustration}
\label{sup:three}

At $N=3$ the single bystander $c$ is also the composite $C$ of \cref{M-prop:composite}. With $\nu=(\nu_B,\nu_c)$, \cref{M-eq:Jsym,M-eq:disturbance} give the unilateral-tariff system
\begin{equation}
J=-\frac{mM_3}{2}\begin{pmatrix}2&-1\\-1&2\end{pmatrix},\qquad
F=-\frac{m}{4}\begin{pmatrix}\rho+\rho_D\\\rho_D-\rho\end{pmatrix},\qquad
M_3=\frac{3\alpha_D(\eta-1)+\rho+1}{2}.
\label{eq:threecountry}
\end{equation}
With $M_3>0$,
\begin{equation}
\frac{d\log e_{AB}}{d\tau}=-\frac{\rho+3\rho_D}{6M_3}<0,\qquad
\frac{d\log e_{Ac}}{d\tau}=\frac{\rho-3\rho_D}{6M_3},\qquad
\frac{d\log e_{Bc}}{d\tau}=\frac{\rho}{3M_3}>0.
\label{eq:threecountryslopes}
\end{equation}
Since $(-J)^{-1}=\frac{2}{3mM_3}\begin{psmallmatrix}2&1\\1&2\end{psmallmatrix}$ is entrywise positive under MML, \cref{M-eq:signrule} gives $d\nu_c/d\tau=\frac{2}{3mM_3}(F_B+2F_c)$, with $F_c=m(\rho-\rho_D)/4$ and $F_B=-m(\rho+\rho_D)/4<0$. For $\rho_D<\rho<3\rho_D$, $F_c>0$ but $F_B+2F_c=m(\rho-3\rho_D)/4<0$: $c$'s trade balance improves at the initial exchange rates, yet $A$ appreciates against $c$ (\cref{M-fig:geometry-a}).

\section{Composite bystander with a fixed internal elasticity}
\label{sup:composite}

If the composite bystander of \cref{M-prop:composite} keeps the mapped $y_C$, $h_C$, flows, weights and $\rho$ but replaces $\eta_C$ by $\eta$, it no longer reproduces the $N$-country economy for $N\geq4$, and its bystander threshold, computed from the linear system, falls slightly with $N$: at the US-like calibration, from $3.96$ at $N=3$ to $3.71$ at $N=8$, while $N\rho_D$ rises from $3.96$ to $10.56$. The steep $N$-dependence of $N\rho_D$ therefore comes from the bloc's internal trade, governed by $\rho$, rather than from its size.

\section{Common and finite tariffs: numerical results}
\label{sup:tau0}

This section reports the numerical values behind \cref{M-app:asym-tariff}. In the reported exercises a common external tariff raises the threshold $\rho^*(\tau_0)$ of \cref{M-eq:tau0thresh}. At $N=3$, $\rho^*(\tau_0)$ is about 10 percent above its free-trade value $N\rho_D$ at $\tau_0=0.25$ and about 36 percent above it at $\tau_0=1$. The main reason is that $x$ falls: the common appreciation of $A$'s currency lowers $A$'s share of each partner's imports, so relative treatment disturbs a smaller part of the trade on which relative exchange-rate movements act. The local threshold applies to marginal cuts; in the exact model, removing the tariff on $B$ entirely requires a higher $\rho$ than a marginal cut at the same baseline.

A concentrated finite tariff works in the opposite direction. At $N=3$, starting from the symmetric free-trade benchmark and re-solving the equilibrium at each tariff, the unilateral bystander threshold falls with the size of the tariff in the computed cases: at the 2025 peak rate $\tau=1.45$ it is $2.70$, against $3\rho_D=3.96$. A larger tariff shrinks $B$'s share of $A$'s imports, so a further increase adds less average protection while continuing to divert expenditure toward $C$.

\section{Asymmetric primitives: further numerical examples}
\label{sup:asym}

Once primitives rather than shares are held fixed, the shares themselves depend on $\rho$: at an asymmetric baseline relative prices differ across origins, and a higher $\rho$ reallocates expenditure toward the cheapest origin. A large target, whose currency depreciates to balance its trade, can then crowd $C$ out of $A$'s imports. With all other primitives symmetric and the baseline re-solved at each $\rho\in[1.02,1000]$, $e_{AC}$ reverses only for $5.99<\rho<53.8$ when $L_B=2.2$ and nowhere when $L_B=2.6$. Away from symmetry there may therefore be no bilateral threshold, one or several. Near the symmetric benchmark, the bilateral threshold rises with $A$'s origin weight on the target (a larger average-protection shock), with the target's size (which raises its share in $A$'s imports) and with home bias in $A$ alone (stronger domestic diversion), and is nearly insensitive to the bystander's size.

The symmetric ordering $d\nu_B<d\nu_C$ can fail away from symmetry: with country sizes $L_A=5.12$, $L_B=4.26$ and $L_C=2.28$, $\beta_{AB}=0.86$, $\beta_{BA}=0.65$, $\beta_{CA}=0.19$, $\alpha_D=0.89$, $\alpha_T=0.6$, $\eta=3.74$ and $\rho=1.20$, the difference $d\nu_B-d\nu_C$ equals $0.08\,d\tau$.

In the example of \cref{M-app:asym-primitives} with distinct bilateral, aggregate terms-of-trade and NEER thresholds, the bystander threshold lies below $3\rho_D$ because the target is small and lightly weighted. A larger bystander weight $w_{AC}$, obtained with a smaller target, lowers all three thresholds.

\section{Illustrative U.S.--China--rest-of-world mapping}
\label{sup:usmap}

This mapping places a unilateral US tariff on China in the three-country model, with $A$ the United States, $B$ China and $C$ the rest of the world. It maps trade shares and country sizes; it is not a calibration of exchange-rate magnitudes.

All inputs are for 2024, in billions of current dollars. US goods imports are $3{,}266.4$ in total and $440.3$ from China, and US goods exports $2{,}061.7$ in total and $143.3$ to China (US Census Bureau). China's goods imports are $2{,}585.1$ in total and $163.6$ from the United States, and its exports $3{,}577.2$ in total and $524.7$ to the United States (China's General Administration of Customs). This gives $b_{AB}=0.135$ and $b_{BA}=0.063$. The rest of the world's imports from the United States are US exports less US exports to China, and its imports from China are China's exports less China's exports to the United States, which gives $b_{CA}=0.386$ and $b_{CB}=0.614$. Country sizes are GDP at purchasing-power parity (IMF World Economic Outlook, October 2026), with the rest of the world equal to world GDP less that of the United States and China; relative to the United States they are $(1,1.30,4.45)$. The origin weights $\beta_{ij}$ are chosen so that the balanced free-trade baseline, solved at $\bar\rho=3\rho_D=3.96$ and then held fixed as $\rho$ varies, reproduces these six shares exactly. The remaining parameters are $\alpha_T=0.4$ and $\eta=1.5$, with either a common $\alpha_D=0.8$ or, for the United States, $\alpha_D=0.673$, the domestic share of US goods absorption in 2024 (BEA goods output plus imports less exports).

The bilateral bystander threshold, above which the dollar depreciates against the rest of the world, is $4.01$ with the common home bias and $3.87$ with the measured US home share. The aggregate terms-of-trade condition \cref{M-eq:totcondition} fails only for $\rho$ above $4.66$ and $4.50$, respectively. MML holds at both baselines for all $\rho\in[1,20]$. The mapping imposes balanced trade on every country, whereas the observed US goods deficit in 2024 was about $1{,}205$; it ignores the gaps between mirror trade statistics, uses model-implied rather than measured home shares for China and the rest of the world, covers goods only and treats the rest of the world as a single country.

Two features of the US--China configuration push in opposite directions. The rapid reallocation of trade between China and other suppliers in the earlier US--China trade war \citep{fajgelbaum2024,freund2024,alfaro2023} suggests high cross-origin substitutability in trade flows, which pushes toward dollar weakness against bystanders. China's size and import weight push against the depreciation mechanism: a larger and more heavily weighted target raises the bystander threshold and moves the aggregate terms-of-trade condition toward its conventional sign. In this mapping, the bilateral bystander threshold for a US tariff on China is about $3.9$--$4.0$ and the aggregate terms-of-trade threshold about $4.5$--$4.7$. A quantitative network model is needed to determine the net sign and the magnitudes. As an interpretation outside the static model, markets may price the tariff configuration expected to persist, in which case expected retaliation would push the dollar in the same direction as realized retaliation.

\section{The 2025 tariff chronology and classification details}
\label{sup:episode}

\paragraph{Tariff chronology.} Statutory headline tariffs above the pre-2025 baseline, by effective date (2025), in percent: US on China 10 (February~4), 20 (March~4), 30 (April~5), 104 (April~9), 145 (April~10), 30 (May~14), 20 (November~10); China on US, across the board, 84 (April~10), 125 (April~12), 10 (May~14), with the product-specific measures of February~10 and March~10 not shown in \cref{M-fig:episode}; US on EU 10 (April~5), 20 (April~9 only), 10 (April~10), 15 inclusive of MFN duties (August~7); US baseline on unlisted partners 10 (April~5); US on Canada 25 (March~4), non-USMCA goods only from March~7, 35 (August~1), with the rate on Mexico at 25 from March~4 to year-end (not shown); Canada on US 25 on covered goods (March~4), with coverage widened on March~13 and April~9 and narrowed to steel, aluminium and autos on September~1.

\paragraph{Product coverage.} China's Phase~1 responses were product-specific tariffs of 10--15 percent, on coal, liquefied natural gas, crude oil and some machinery and vehicles from February~10 and on agricultural products from March~10.

\paragraph{Canada and Mexico.} The US tariffs on Canada (10 percent on energy) and Mexico took effect on March~4 and from March~7 no longer applied to goods qualifying under the USMCA. Canada retaliated on March~4 with a 25 percent surtax on C\$30 billion of US goods \citep{canadafinance2025} and extended it to a further C\$29.8 billion on March~13; Mexico did not retaliate. Canada and Mexico together carry 27.6 percent of the broad index. In the week of March~3--7 the broad index fell by $1.8$ percent, more than its net change over Phase~1, and the dollar fell by $4.4$ percent against the euro.

\paragraph{Germany's fiscal announcement.} The week of March~3--7 also contained the March~4 announcement by Germany's prospective governing parties of a large defence and infrastructure package and a loosening of the debt brake \citep{ripper2025}.

\paragraph{Vietnam.} The dollar strengthened against the dong by $3.2$ percent between the January and December averages, all of it by July. Vietnam, however, moved during the year from a bystander-like position to a direct target: a 46 percent rate was announced on April~2 and was in force for one day before the suspension of April~10, and a 20 percent rate applied from August~7. From August~7 to November~10, Vietnam's 20 percent against China's 30 percent was relative treatment in Vietnam's favour; after November~10, when both faced 20 percent, it disappeared. Within Phase~3 the dollar was roughly flat against the dong ($0.0$ percent on monthly averages, $+0.5$ percent between end-of-month rates).

\paragraph{Movements outside the phases.} Between the May~12 and August~6 daily observations, which fall outside the phases, the dollar fell by a further $4.9$ percent against the euro and $2.5$ percent against the broad index; in December it was $13.1$ percent weaker against the euro than in January.

\section{Sector-specific factors}
\label{sup:rv}

Replace the linear tradable technology with a Ricardo--Viner technology, $Y_{Ti}=A_{Ti}L_{Ti}^{\gamma}$ with $0<\gamma\leq1$, where a fixed sector-specific factor owned by domestic households earns rents $(1-\gamma)Y_{Ti}$; non-tradables remain linear in labor, which is mobile across sectors. At a free-trade baseline, tradable employment is $L_{Ti}=\gamma\alpha_TL_i/(1-\alpha_T+\gamma\alpha_T)$ whatever the exchange rates, so the adjustment matrix is that of the linear model with the same nominal incomes. Only the disturbance changes, through the tariff-setter's income: the rebate draws labor out of tradables and raises the wage, so with $s_{AB}$ the share of $A$'s income spent on $B$'s good,
\begin{equation*}
\frac{d\log I_A}{d\tau_{AB}}\bigg|_{\nu}=\kappa(\gamma)\,s_{AB},\qquad
\kappa(\gamma)\equiv1+\frac{(1-\alpha_T)^2(1-\gamma)}{\alpha_T\,(1-\alpha_T+\gamma\alpha_T)}\geq1,
\end{equation*}
with equality if and only if $\gamma=1$. \Cref{M-prop:unilateral} then holds with $\rho_D$ replaced by
\begin{equation*}
\rho_D(\gamma)\equiv\rho_D-(\kappa(\gamma)-1)m=\alpha_D\eta+(1-\alpha_D)(1-\alpha_T)\,\frac{\gamma}{1-\alpha_T+\gamma\alpha_T}\leq\rho_D,
\end{equation*}
so the bystander threshold falls to $N\rho_D(\gamma)$, while the NEER still appreciates because $\rho_D(\gamma)>0$. At the US-like parameters and $N=3$, the threshold is $3.96$ with linear technology and $3.83$ at $\gamma=1/2$.

\section{Permanent trade imbalances}
\label{sup:imbalance}

Let $T_i$ be permanent transfers fixed in $A$'s currency, positive for a receipt, with $\sum_{i\in\mathcal I}T_i=0$. Each country's external-balance condition becomes
\begin{equation*}
\TB_i+T_i=0\qquad\text{for every }i\in\mathcal I,
\end{equation*}
so a recipient runs a trade deficit equal to its receipt. The general linear system \cref{M-eq:response} remains valid at any such baseline.

As an illustration of U.S. scale, let $A$ receive a transfer equal to 3 percent of its income, paid by the target $B$, with the US-like parameters and $N=3$. Because $A$'s import share is small, financing the deficit requires a baseline appreciation of $A$ of about $15$ log points against $B$. The nominal bystander threshold rises from $3.96$ to $4.94$: $A$'s imports tilt toward the cheapened target, so the bystander is less exposed to diversion from $A$'s tariff. The single-threshold structure survives on $\rho\in[1.05,60]$. Alternative transfer allocations and tariff levels are reported in the replication package.

\clearpage
\singlespacing
\bibliographystyle{aer}
\bibliography{refs}

\end{document}
